\documentclass[11pt,a4paper]{article}

\usepackage[french]{babel}
\usepackage[T1]{fontenc}
\usepackage[utf8]{inputenc}
\usepackage{amsmath,amssymb,amsfonts}
\usepackage{bm}
\usepackage{geometry}
\usepackage{hyperref}

\geometry{margin=2.5cm}

\title{Calcul de la charge topologique sur un maillage triangulaire\\
selon la formulation de Berg et Lüscher}
\author{}
\date{}

\begin{document}

\maketitle

\section{Introduction}

En micromagnétisme, la charge topologique d'un champ d'aimantation unitaire
$\mathbf{m}:\Omega\rightarrow S^2$ est définie dans le cas continu par

\begin{equation}
Q=\frac{1}{4\pi}
\int_{\Omega}
\mathbf{m}\cdot
\left(
\frac{\partial\mathbf{m}}{\partial x}
\times
\frac{\partial\mathbf{m}}{\partial y}
\right)
\,dx\,dy.
\end{equation}

La quantité

\[
\rho=
\frac{1}{4\pi}
\mathbf{m}\cdot
(\partial_x\mathbf{m}\times\partial_y\mathbf{m})
\]

est appelée densité topologique.

Lorsque le domaine est discrétisé par un maillage d'éléments finis, le champ
n'est connu qu'aux sommets des triangles. Les dérivées spatiales deviennent
alors dépendantes de l'interpolation choisie et la formule précédente ne
préserve plus exactement le caractère topologique de la charge.

La méthode proposée par Berg et Lüscher consiste à remplacer l'intégrale locale
par l'aire orientée du triangle sphérique défini par les trois vecteurs
d'aimantation situés aux sommets de chaque élément.

\section{Triangle sphérique}

Considérons un triangle du maillage dont les sommets portent les vecteurs
unitaires

\[
\mathbf{a},\qquad
\mathbf{b},\qquad
\mathbf{c}.
\]

Ces trois vecteurs définissent trois points sur la sphère unité et délimitent
un triangle sphérique.

La contribution topologique de cet élément est définie par

\[
Q_T=\frac{\Omega}{4\pi},
\]

où $\Omega$ désigne l'aire orientée du triangle sphérique.

\section{Aire sphérique}

L'aire d'un triangle sphérique est donnée par son excès sphérique

\[
\Omega=A+B+C-\pi,
\]

où $A$, $B$ et $C$ sont les angles du triangle sphérique.

Cette expression est peu adaptée aux calculs numériques puisque les angles
deviennent difficiles à évaluer lorsque le triangle est petit.

Berg et Lüscher utilisent une formulation entièrement équivalente fondée sur
l'angle solide associé aux trois vecteurs.

\section{Formule de Berg et Lüscher}

On définit

\begin{equation}
N=
\mathbf{a}\cdot
(\mathbf{b}\times\mathbf{c}),
\end{equation}

et

\begin{equation}
D=
1+
\mathbf{a}\cdot\mathbf{b}
+\mathbf{b}\cdot\mathbf{c}
+\mathbf{c}\cdot\mathbf{a}.
\end{equation}

On montre alors que

\begin{equation}
\tan\left(\frac{\Omega}{2}\right)
=
\frac{N}{D}.
\end{equation}

Afin de choisir automatiquement le bon quadrant, on écrit

\begin{equation}
\boxed{
\Omega
=
2\operatorname{atan2}(N,D)
}
\end{equation}

avec

\[
N=
\mathbf{a}\cdot(\mathbf{b}\times\mathbf{c}),
\]

et

\[
D=
1+\mathbf{a}\cdot\mathbf{b}
+\mathbf{b}\cdot\mathbf{c}
+\mathbf{c}\cdot\mathbf{a}.
\]

Cette formule fournit directement l'aire orientée du triangle sphérique.

\section{Charge topologique discrète}

Chaque triangle contribue donc à la charge totale selon

\begin{equation}
\boxed{
Q_T=
\frac{1}{2\pi}
\operatorname{atan2}
\left(
\mathbf{a}\cdot(\mathbf{b}\times\mathbf{c}),
1+\mathbf{a}\cdot\mathbf{b}
+\mathbf{b}\cdot\mathbf{c}
+\mathbf{c}\cdot\mathbf{a}
\right)
}
\end{equation}

La charge totale est simplement obtenue par sommation

\begin{equation}
Q=
\sum_T Q_T.
\end{equation}

\section{Lien avec la définition continue}

Supposons que le triangle soit suffisamment petit.

On écrit

\[
\mathbf{b}
=
\mathbf{a}+\delta_x,
\qquad
\mathbf{c}
=
\mathbf{a}+\delta_y,
\]

avec

\[
\|\delta_x\|,
\|\delta_y\|
\ll1.
\]

Comme les vecteurs sont unitaires,

\[
\mathbf{a}\cdot\delta_x
=
\mathbf{a}\cdot\delta_y
=
0.
\]

Le numérateur devient

\[
N
=
\mathbf{a}\cdot
(\delta_x\times\delta_y)
+O(h^3),
\]

tandis que

\[
D
=
4+O(h^2).
\]

Il vient alors

\[
\Omega
=
2\arctan\left(\frac{N}{4}\right)
=
\frac12
\mathbf{a}\cdot
(\delta_x\times\delta_y)
+O(h^3).
\]

En identifiant

\[
\delta_x
=
\frac{\partial\mathbf{m}}{\partial x}\Delta x,
\qquad
\delta_y
=
\frac{\partial\mathbf{m}}{\partial y}\Delta y,
\]

et en remarquant que l'aire du triangle plan vaut

\[
A_T
=
\frac12
\Delta x\,\Delta y,
\]

on obtient

\[
\Omega
=
A_T\,
\mathbf{m}\cdot
\left(
\frac{\partial\mathbf{m}}{\partial x}
\times
\frac{\partial\mathbf{m}}{\partial y}
\right)
+O(h^3).
\]

La contribution discrète devient alors

\[
Q_T
=
\frac{1}{4\pi}
\int_T
\mathbf{m}\cdot
\left(
\partial_x\mathbf{m}
\times
\partial_y\mathbf{m}
\right)
\,dS
+O(h^3),
\]

ce qui montre que la formulation de Berg et Lüscher est consistante avec la
définition continue.

\section{Propriétés}

La méthode possède plusieurs propriétés importantes :

\begin{itemize}
\item elle ne dépend pas de la forme géométrique des triangles ;
\item elle ne nécessite aucune approximation des dérivées de
l'aimantation ;
\item elle est robuste même lorsque l'aimantation varie fortement ;
\item elle converge vers la définition continue lorsque le maillage est
raffiné ;
\item elle préserve beaucoup mieux le caractère entier de la charge
topologique que les méthodes basées sur les dérivées numériques.
\end{itemize}

\section*{Référence}

B. Berg and M. Lüscher,
\emph{Definition and Statistical Distributions of a Topological Number in the
Lattice O(3) Sigma Model},
Nuclear Physics B, vol. 190, pp. 412--424, 1981.

\end{document}